\cotefile{10}
\scanpage{102}

Quelques exemples de catégories tensorielles.

1) Soit $M$ un groupe. Soit $\underline{C}_M$ la catégorie des vectoriels de dimension finie sur le corps $k$, munis d'une graduation de type $M$. C'est une catégorie tensorielle sur $k$, munie d'un foncteur fibre sur $k$, le foncteur ``\unsure{oublif|oublie|oubli} de la graduation''. Le groupe algébrique associé est le groupe de type multiplicatif $D_k(M)$ (SGA 3 I 4.7.3). Par exemple si $M=\underline{Z}^r$, on trouve $G=\underline{G}_m^r$.

Application: soit $\underline{C}$ une catégorie tensorielle sur $k$ munie d'un foncteur fibre \inserted{$F$} sur $k$, donc associée à un schéma en groupes affine $G$ sur $k$. On cherche toutes les façons de mettre, pour chaque $M\handwritten{\in}\operatorname{Ob}\underline{C}$, une graduation de type $M$ sur $F(V)$, de façon fonctorielle en $M$, et compatible (dans un sens évident) avec les produits tensoriels.Elles correspondent aux $\unsure{\otimes|\Omega}$-foncteurs de $\underline{C}$ dans $\underline{C}_M$ compatibles avec les foncteurs fibres, donc aux homomorphismes de $D_k(M)$ dans $G$. Par exemple, si $M=\underline{Z}^r$, il faut prendre les homomorphismes de $\underline{\unsure{G|E|\mathbb{G}}}_m^r$ dans $G$.

Dans la situation précédente, on peut se demander quand une $\otimes$-graduation de type $M$ du foncteur $F$ correspond à une graduation de type $M$ du foncteur identique de $\underline{C}$, i.e. pour tout $V$, la graduation de $F(V)$ provient d'une graduation de $V$. On trouve qu'il faut et il suffit pour cela que l'homomorphisme correspondant $D_k(M)\handwritten{\longrightarrow}G$ soit central. Par exemple, si $\underline{C}$ est la catégorie des motifs sur $k$, et si on dispose d'un foncteur fibre $F$ de $\underline{C}$ sur $k$, alors on trouve un homomorphisme central canonique $i:\underline{G}_m\handwritten{\longrightarrow}G$. D'ailleurs, la donnée du motif de Tate (qui est de rang $2$, \unsure{èt|et} de ``poids'' $2$ pour la graduation naturelle) correspond à \deleted{un} la donnée d'un homomorphisme $j:G\handwritten{\longrightarrow}\underline{G}_m=\operatorname{Gl}(1)$. Le fait que $T$ soit de poids $2$ s'exprime par la relation
\[
ji(\handwritten{\lambda})=\handwritten{\lambda}^{2}.
\]
% archive sheet number: 101
\scanpage{103}

Lorsque $k$ est de caractéristique nulle, on a toujours un foncteur fibre naturel: le \uline{foncteur de Hodge,} qui à la cohomologie motivique d'une variété (projective lisse) $X$ associe le vectoriel bigradué $\handwritten{\coprod}H^q(X,\handwritten{\underline{\Omega}}^p_{X/k})$ Donc pour le groupe de Galois motivique correspondant $G$, on trouve un homomorphisme naturel
\[
\underline{G}_m^2\handwritten{\longrightarrow}G,
\]
i.e. deux homomorphismes commutant l'un à l'autre
\[
i_1,i_2:\underline{G}_m\handwritten{\longrightarrow}G.
\]
Le fait que la graduation totale dans la cohomologie de Hodge corresponde au poids des motifs s'exprime par la relation
\[
i(\handwritten{\lambda})=i_1(\handwritten{\lambda})\,i_2(\handwritten{\lambda});
\]
on fera attention que $i_1$ et $i_2$ ne sont pas centraux (car la bigraduation en cohomologie de Hodge ne correspond pas à une bigraduation d'un motif !)

2) Prenant toujours pour $\underline{C}$ la catégorie des motifs sur $k$, avec $k$ de car. nulle, on a un autre foncteur fibre canonique, le \uline{foncteur de De Rham} qui associe à la cohomologie motivique d'une variété $X$ le vectoriel $H^*(X,\handwritten{\Omega}^*_{X/k})$ (espace d'hypercohomologie). Ce vectoriel n'est plus bigradué\handwritten{\unsure{é,|en,}} mais seulement gradué et filtré,\handwritten{(}la filtration étant celle associée à la suite spectrale d'hypercohomologie, commençant avec la cohomologie de Hodge\handwritten{)}. On sait d'ailleurs que cette suite spectrale dégénère (théorie de Hodge), donc $\operatorname{Gr}(H_{\mathrm{DR}}(V))\handwritten{\simeq}H_{\mathrm{Hdg}}(V)$ (isomorphisme fonctoriel en le motif $V$\inserted{\handwritten{où on \unsure{gradue} le $2^{\mathrm{ème}}$ membre \uline{par $p$}}}). On peut se proposer d'analyser à quelle structure supplémentaire, sur le groupe de Galois motivique associé au foncteur fibre $H_{\mathrm{DR}}$, correspond la filtration canonique de ce foncteur.
% The inserted handwritten note is circled, with a pointer immediately after V.

De façon générale, étant donné une catégorie tensorielle $C$ sur $k$ munie d'un foncteur fibre $F$, on peut se proposer de déterminer les filtra%
% archive sheet number: 102
\scanpage{104}%
tions sur $F$ (décroissantes, \inserted{discrètes,} indexées par $\underline{Z}$) compatibles avec lesproduits tensoriels (en utilisant la notion év\deleted{\unsure{o|q}}\handwritten{i}dente de produit tensoriel\deleted{s} de deux espaces vectoriels filtrés). On notera qu'une telle donnée ne pourra plus s'exprimer par un homomorphisme d'un certain groupe algébrique dans $G$ (le groupe de Galois de $\underline{C}$ en $F$), car la catégorie des espaces vectoriels de dimension finie sur $k$, munis d'une filtration décroissante discrète indexée par $\underline{Z}$, n'est pas une catégorie abélienne (les bimorphismes ne sont pas des isomorphismes\handwritten{!}). Mais à une telle donnée est \unsure{associée|associé} un deuxième foncteur fibre\margin{\handwritten{$F'$}} $\inserted{\unsure{F|\mathbb{F}}'(V)=}\operatorname{Gr}F(V)$, à valeurs cette fois-ci dans les vectoriels gradués (NB $F$ joue le rôle de la cohomologie de \deleted{Hodge, \unsure{$F'$}} De Rham, $F'$ celui de la cohomologie de Hodge\deleted{\unsure{)}}, muni de la graduation \deleted{totale}\inserted{\handwritten{par $p$}}). La donnée de $F'$ correspond alors à un torseur \inserted{\handwritten{(à droite)}} $P$ sous $G$, le groupe de Galois correspondant $G'$ étant déduit de $G$ en tordant par $P$. La graduation de $F'$ correspond à un homomorphisme $i_{\handwritten{1}}:\underline{G}_m\handwritten{\longrightarrow}G'$. Considérons alors sur $F'$ la filtration décroissante associée à sa graduation, et soit $H'_{i_{\handwritten{1}}}=\underline{\operatorname{Aut}}\,\underline{\mathrm{filt}}^{\handwritten{!}}(F')\handwritten{\subset}\underline{\operatorname{Aut}}(F')=G'$ le sous-schéma en groupes de $G$ qui correspond aux \inserted{$\unsure{\otimes|\Omega}$-}automorphismes de $F'$ (ou plutôt des $F'_{k'}$, $k'$ une $k$-algèbre quelconque) qui respectent sa filtration\inserted{\handwritten{et induisant l'identité sur le gradué associé}}; il est canoniquement déterminé par $i_{\handwritten{1}}$. On peut aussi regarder le sous-schéma
\[
Q=\underline{\operatorname{Isom}}\,\underline{\mathrm{filt}}^{\handwritten{!}}(F,F')\handwritten{\subset}P=\underline{\operatorname{Isom}}(F,F')
\]
du schéma $P$ des $\otimes$-isomorphismes de $F$ avec $F'$, \deleted{formé de} qui correspond aux automorphismes respectant les filtrations de $F$ et $F'$\inserted{\handwritten{et induisant l'identité sur les gradués associés}}. C'est à priori un pseudo-torseur à \deleted{droite}\inserted{\handwritten{gauche}} sous $H'$, i.e. il est vide ou un torseur à droite sous $H'$. \uuline{Il faudrait prouver} que c'est bien un torseur (i.e. que sur une extension convenable $k'$ de $k$, on peut trouver un isomorphisme de $F_{k'}$ avec $F'_{k'}$ \uline{respectant les filtrations}). Donc on trouve une restriction du groupe d'opérateurs (à \deleted{droite}\inserted{\handwritten{gauche}}) $G'$ de $P$ au sous-groupe $H'$,
% archive sheet number: 103
\scanpage{105}%
par un $H'$-torseur à gauche $Q$.

Moyennant la vérification laissée en suspens à l'instant, on trouve alors que la donnée d'un ``foncteur fibre \uline{filtré}'' $F$ de $\underline{C}$ sur $k$ revient à la donnée

(i) d'un foncteur fibre $F'$ \deleted{gradué de type} (d'où un \unsure{grouped|groupe} de Galois $G'=\underline{\operatorname{Aut}}(F')$ );

(ii) D'une graduation de type $\underline{Z}$ de $F'$, i.e. un homomorphisme
\[
i_1:\underline{G}_m\handwritten{\longrightarrow}G',
\]
ce qui permet de définir comme ci-dessus un sous-groupe $H'_{i_1}$ de $G'$ ;

(iii) D'un torseur à gauche $Q'$ sous $H'_{i_1}$.

A ces données, on associe simplement le foncteur fibre tordu
\[
F=F'\mathop{\handwritten{\wedge}}\limits^{H'_{i_1}}Q',
\]
$F$ étant filtré par la filtration déduite de celle de $F'$ \deleted{(\unsure{via} la \unsure{graduation|graduatio}} (associée à la graduation de $F'$) en tordant par $Q'$.

\deleted{NB Dans le cas où $F$ est le foncteur de De Rham sur la catégorie des motifs sur $k$, on prouve (moyennant le yoga général des motifs - il est possible qu'on utilise pour ceci les conjectures de Hodge ...) que le groupe $H'$ est une limite projective de groupes algébriques \deleted{\inserted{\unsure{lisses}}} \deleted{\unsure{à composantes} connexes \uline{résolubles}\unsure{.} \unsure{Donc}}\inserted{unipotents} Si on se borne à une catégorie de motifs \deleted{\unsure{de|en}} à engendrement fini, et si $G$ est connexe, on trouve aussi que $H'$ est connexe, donc étant résoluble, on trouve}

On constate aisément que le groupe $H'$ est nécessairement une limite projective de groupes algébriques \uline{unipotents}. On en conclut aussitôt que si $\underline{C}$ est à engendrement fini (ou, plus généralement, \inserted{\handwritten{à \unsure{engendrement}}} dénombrable, de façon que $G'$ \deleted{soit} donc aussi $H'$ soit limite projective d'une \uline{suite} de groupes algébriques) alors tout torseur sous $H'$ est trivial; cela signifie ici que tout foncteur fibre filtré est en fait associé à un foncteur fibre gradué (en prenant la filtration correspondant à la graduation)
% archive sheet number: 104
\scanpage{106}%
i.e. que la filtration dudit foncteur admet un\deleted{e} splittage compatible avec les produits tensoriels. Mais le choix d'un tel splittage équivaut à celui d'un point \deleted{sous} sur un certain torseur à droite $Q=Q'$ sous le \inserted{schéma en} groupes \inserted{(pro-unipotent)} $H$ des automorphismes de $F$ respectant la filtration et induisant l'identité sur le gradué associé; il n'est pas du tout canonique !

3) Appelons \uline{pré-structure de Hodge} sur un espace vectoriel \inserted{$V$} de dimension finie sur $\underline{Q}$, la donnée d'une bigraduation sur $V\unsure{\otimes|\Omega}_{\underline{Q}}\underline{C}=V_{\underline{C}}$, $V_{\underline{C}}=\handwritten{\coprod_{p,q}}V^{p,q}$, telle que a) la graduation totale correspondante soit ``définie sur $\underline{Q}$'' \uline{i.e.} $\handwritten{\coprod}_{p+q=n}V^{p,q}$ provienne d'un sous espace $V_{\underline{Q}}^n$ de $V$, et b) on a $\overset{\handwritten{\rule{.65em}{.35pt}}}{V}^{p,q}=V^{q,p}$, où $x\handwritten{\mapsto}\overset{\handwritten{\rule{.55em}{.35pt}}}{x}$ désigne la conjugaison complexe. \unsure{(}NB Généralisation à des corps plus généraux laissée à Saavedra). Les vectoriels $V$ munis d'une pré-structure de Hodge forment une catégorie tensorielle sur $\underline{Q}$ dans un sens évident, muni d'un foncteur fibre canonique, le foncteur ``oubli''\inserted{$F$}. On trouve donc un groupe de Galois $G$, et \deleted{tout} plus généralement toute \inserted{$\otimes$-}sous-catégorie $\underline{C}$ de la catégorie précédente nous définit un groupe de Galois $G$. La bigraduation sur le foncteur \deleted{oubli} $F_{\underline{C}}(V)=V\otimes_{\underline{Q}}\underline{C}$ correspond, en vertu de 1) (où il convient cependant de se permettre une extension du corps de base sur le foncteur fibre envisagé ....) d'un homomorphisme ${\underline{G}_m}_{\underline{C}}^2\handwritten{\longrightarrow}G_{\underline{C}}$, i.e. de deux homomorphismes \inserted{qui commutent}
\[
i_1,i_2:{\underline{G}_m}_{\underline{C}}\handwritten{\longrightarrow}G_{\underline{C}}.
\]
Les deux conditions a) et b) imposées aux structures envisagées s'interprètent respectivement par les faits que l'homomorphisme
\[
i_{\underline{C}}=i_1i_2\unsure{\doteqdot|\div|=}\handwritten{\lambda\mapsto}i_1(\handwritten{\lambda})i_2(\handwritten{\lambda})
\]
est ``défini sur $\underline{Q}$'' i.e. provient d'un homomorphisme
\[
i:\underline{G}_m\handwritten{\longrightarrow}G,
\]
(nécessairement central, car les composantes homogènes $V_{\underline{Q}}^n$ d'une pré-structure de Hodge sont \deleted{nécessairement}\inserted{évidemment} munis d'une \inserted{pré-}structure de Hodge de façon que $V$ soit la somme directe de $V_{\underline{Q}}^n$ en tant que pré-structure
% archive sheet number: 105
\scanpage{107}%
de Hodge), et par la condition que l'on a
\[
i_2=\overset{\handwritten{\rule{0.8em}{0.4pt}}}{i_1}\quad\text{i.e.}\quad i_2(\handwritten{\lambda})=\overset{\handwritten{\rule{2.2em}{0.4pt}}}{i_1(\handwritten{\bar{\lambda}})}\quad\text{pour tout}\quad\handwritten{\lambda\in}\underline{C}.
\]

Si à toute variété projective lisse $X$ sur $\underline{C}$ on associe sa cohomologie rationnelle $H^{\unsure{*|'}}(X,\underline{Q})=V$, de sorte que $V_{\underline{C}}=H^*(X,\underline{C})$ est isomorphe canoniquement (par la théorie de Hodge) à $H_{\mathrm{Hdg}}(X)=\handwritten{\coprod}H^q(X,\handwritten{\underline{\Omega}}^p_{X/\underline{C}})$, on voit qu'on trouve ainsi une pré-structure de Hodge sur $H^*(X,\underline{Q})$, d'où un \inserted{$\unsure{\otimes|\Omega}$-}foncteur naturel de la catégorie des motifs sur $\underline{C}$ dans la catégorie des pré-structures de Hodge. La conjecture de Hodge \unsure{stándart|standart} équivaut à dire que ce foncteur est \uline{pleinement fidèle}, i.e. que l\deleted{e}'homomorphisme naturel qui va du groupe de Galois de Hodge précédent $G$ dans le groupe de Galois motivique (associé au \inserted{$\otimes$-}foncteur de Betti $H_{\mathrm{Bet}}$) est un épimorphisme. \unsure{OU|Ou} encore que pour toute catégorie \inserted{$\underline{C}_o$} \inserted{de type fini} de motifs sur $\underline{C}$, désignant par $\underline{C}$ la \inserted{$\otimes$-}catégorie de pré-structures de Hodge engendrée par \deleted{\unsure{celle-ci,}} les $H_{\mathrm{Bet}}(M)$ pour $M\handwritten{\in}\operatorname{Ob}(\underline{C}_o)$, l'homomorphisme de groupes algébriques $G\handwritten{\longrightarrow}G_o$ associé au foncteur de Betti-Hodge $\underline{C}_o\handwritten{\longrightarrow}\underline{C}$ est un épimorphisme (i.e. surjectif sur les points à valeurs complexes, disons).
% A curved transposition mark takes “de type fini” from after “de motifs” to the caret before it.

On appelle \deleted{structure} \uline{polarisation} d'une pré-structure de Hodge de poids $n$ la donnée d'un accouplement de préstructures de Hodge
\[
\handwritten{\phi}:V\mathbin{\mathrm{x}}V\handwritten{\longrightarrow}\underline{Q}(n),
\]
où $\underline{Q}(n)$ est l'espace vectoriel trivial $\underline{Q}$ sur $\underline{Q}$, avec $\underline{Q}_{\underline{C}}=\underline{C}$ muni du \deleted{\unsure{la}} bidegré $(n,n)$, ayant la propriété que la forme hermitienne correspondante sur $V_{\underline{C}}$,
\margin{\handwritten{\uline{$\unsure{\psi|\phi}$}}}
\[
\handwritten{\psi}(x,y)=\handwritten{\phi}(x,\overset{\handwritten{\rule{0.5em}{0.4pt}}}{y})(\handwritten{i})^{p\handwritten{-q}}\quad\text{pour }x\text{ de bidegré }(p,\deleted{\unsure{n}}\handwritten{q}\unsure{-}p)
\]
soit définie positive. Une \uline{structure de Hodge} est une préstructure de Hodge admettant une polarisation. Les structures de Hodge forment une sous-$\otimes$-catégorie de la catégorie des pré-structures de Hodge. La théorie de Hodge nous assure que le foncteur de Betti-Hodge sur la catégorie des
% archive sheet number: 106
\scanpage{108}%
motifs sur $\underline{C}$ prend ses valeurs en fait dans la catégorie des structures de Hodge (une polarisation d'une varéiét projective lisse $V$ définit canoniquement une polarisation de la structure de Hodge associée sur la cohomologie de Betti-Hodge). NB On n'a aucune idée sur ce que pourrait être l'image essentielle du foncteur précédent, par exemple s'il y a lieu d'espèrer qu'on trouve toutes les structures de Hodge (doncune équivalence de catégories : motifs sur $\underline{C}$ $\handwritten{\longrightarrow}$ structures de Hodge); cela semble peu probable, mais on n'a aucune indication sérieuse dans un sens ou l'autre.

La catégorie des \deleted{pré} structures de Hodge est semi-simple. \deleted{En fait,} Si $G$ est le groupe de Galois d'une \inserted{$\unsure{\otimes|\Omega}$-}catégorie \inserted{$\underline{C}$} de pré-structures de Hodge, à engendrement fini si on veut (pour simplifier), alors on peut expliciter en termes du groupe de Galois associé et de s\deleted{\unsure{e}}a\deleted{s} structure\deleted{s} $i_1\deleted{,\unsure{i_2}}$ ci-dessus la condition \deleted{que} pour que les objets de $\underline{C}$ soient en fait des structures de Hodge. Ceci est un exercice plaisant et délectable, qui devrait figurer dans un travail systématique sur les $\otimes$-catégories, dans le chapitre des exemples. On trouve des restrictions très sérieuses sur le groupe $G$ \deleted{(plus} muni de $i_1$ (en plus du fait que $G$ soit réductif).

4) Je laisse le soin à Saavedra de déterminer quelle structure supplémentaire on obtient sur la structure de Hodge ``complexe'' associée à une variété projective \inserted{lisse} complexe $X$, lorsqu'on se donne cette dernière comme déduite d'une variété projective réelle $X_{\underline{R}}=X_o$.On trouve une notion de ``structure de Hodge réelle'', donnant naissance à une $\otimes$-catégorie correspondante. Dans le groupe de Galois motivique de celui-ci, en plus de la structure $i_1$, on trouve un élément $f_{\handwritten{\infty}}$ de $G(\underline{Q})$, \deleted{désigné} d'ordre 2 (jouant le rôle d'un ``élément de Frobenius à l'infini''), qui correspond à l'automorphisme du foncteur de Betti $X_o\handwritten{\longmapsto} H^*(X_o(\underline{C}),\underline{Q})$ déduit de l'homéomorphisme $x\handwritten{\mapsto}\overset{\handwritten{\rule{.55em}{.35pt}}}{x}$ de $X_o(\underline{C})$. Il faut expliciter les relations entre cet élément et $i_1$, $i_2$ !
% archive sheet number: 107
\scanpage{109}

5): Soit $\underline{C}$ une $\unsure{\otimes|\Omega}$-catégorie tensorielle, \deleted{Pour} munie d'un foncteur fibre sur $k$ \inserted{de caractéristique nulle.} A prouver que, pour que le groupe de Galois $G$ correspondant soit profini, il faut et il suffit que pour tout objet $M$ de $\underline{C}$, la $\otimes$-catégorie engendrée soit semi-simple et n'ait qu'un nombre fini d'objets simples non isomorphes. Si $\underline{C}$ est quelconque, la sous-catégorie pleine \deleted{de} $\underline{C}_o$ de $\underline{C}$ qui correspond au pro-groupe quotient de $G$ formé des $\inserted{G_i\handwritten{/}}G_i^o$ (où $G=\underset{\handwritten{\longleftarrow}}{\lim}G_i$, et $G_i^o$ est la composante neutre de $G_i$) est formée exactement des objets $M$ ayant la propriété précédente.

Il serait intéressant de trouver des énoncés correspondants en caractéristique quelconque.

6) La notion de polarisation d'un motif sur un corps (elle-même déduite de celle de polarisation d'une variété algébrique) donne une structure supplémentaire remarquable dans la catégorie des motifs: si $M$ est un motif de poids $n$, on sait parmi les formes symétriques ($n$ pair) resp. alternées ($n$ impair) $M\otimes M\handwritten{\unsure{\rightharpoonup|\longrightarrow}} T(n)$ (où $T$ est le motif de Tate) distinguer celles qui sont ``définies positives'' ou encore des \handwritten{``}polarisations\handwritten{''}. Cette notion se reflête par exemple par des structures supplémentaires sur les groupesde Galois motiviques. Il y a lieu de faire une étude axiomatique abstraite d'une telle notion de polarisation sur une $\otimes$-catégorie générale au dessus d'un \deleted{corps} sous-corps du corps des réels. On pourra en rediscuter à l'occasion.
% archive sheet number: 108
